\section{Introduction}
\label{sec:introduction}

Distributed multi-object tracking combines complementary sensor views without
centralizing measurements. Labeled multi-Bernoulli (LMB) filters jointly
represent target existence, state uncertainty, and identity
\cite{Vo2014LRFS,Reuter2014LMB,Vo2019MSGLMB}. A Gaussian-mixture (GM) LMB
message may contain several components per label. The studied implementation
serializes every component although its receiver immediately collapses each
label to one Gaussian. It therefore transmits fields this receiver cannot
observe.

Communication-aware RFS trackers reduce transmission frequency, select mixture
components, or exchange projected state estimates
\cite{Shen2022EventTriggeredLMB,Li2023EventTriggeredConsensusLMB,
Shen2024IntegralEventLMB,
Li2019PartialConsensusGMPHD,Xue2026StateCovProjection}. These methods decide
when to send, which components to share, or which compact estimate to form. We
instead ask a consumer-facing question: for a fixed receiver, which transmitted
fields can affect its output?

We study a projected Gaussian KLA-LMB receiver with one synchronous neighbor-
fusion round per filtering step. It aligns labels, normalizes Metropolis weights
over sources carrying each label, projects each GM Bernoulli to existence and
two moments, and fuses the resulting Gaussians. Let $\mathcal T$ denote the
canonical encode/decode content map. For the executed receiver
$\mathcal F_{\omega,\mathcal R}=\mathcal G_{\omega,\mathcal R}\circ\mathcal P$,
the wire-path identity is
$\mathcal F_{\omega,\mathcal R}(\mathcal T\pi)=
\mathcal F_{\omega,\mathcal R}(\mathcal T\mathcal P\pi)$. The nontrivial check
is that label presence, Gaussian overlap, and deterministic regularization also
factor through $\mathcal P$ after transport
(Fig.~\ref{fig:receiver-factorization}).

Our contributions are threefold. First, we extract an executable receiver
contract that includes label presence, KLA weights, covariance canonicalization,
and solve regularization. Second, a versioned application-layer codec and
adversarial property tests validate the contract across serialization. Third, a
frozen 50-trial paired evaluation shows a 58.28\% mean attempted-byte reduction
and zero residual across 1,119,037 label-instance comparisons. We claim neither
equality of the underlying mixture densities nor sufficiency for a mixture-aware
or multi-round receiver.
